Named entity recognition (NER) for Chinese equipment-domain documents must
operate with scarce annotations while limiting false-positive entity
predictions from generative large language model (LLM) extractors. This setting
requires control over what evidence enters each prediction and a traceable
record of how that evidence is used. We construct a small frozen dataset from
public equipment-related texts, with 30 support passages, 50 development
passages, and 170 test passages, and propose Evidence-Controlled Few-Shot
Information Extraction (EC-FSIE). EC-FSIE coordinates query-adaptive
demonstration selection, bounded schema-specific terminology, deduplicated and
reranked RAG evidence, and audit-preserving execution around a local
\modelname{Qwen2.5-7B-Instruct} extractor. We compare four progressively
controlled NER configurations on the frozen test set: fixed-random
demonstrations, $k$-nearest-neighbor (kNN) demonstrations, kNN demonstrations
with a domain glossary, and the same configuration augmented with reranked
retrieval-augmented generation (RAG) evidence. A fresh two-human calibration on
20 previously unused non-test passages obtained exact span-plus-type micro-F1
of 0.8000, meeting the study's pre-specified manuscript threshold. The
class-specific audit also identifies \entitytype{TECH\_SCHEME} (F1 0.153846) as
the clearest target for future annotation refinement. In descriptive
three-seed ablations at $K=5$, kNN improved F1 by 0.020524 over the
fixed-random baseline, and glossary injection added 0.013689 F1 mainly through
recall. The glossary-only configuration C and the RAG configuration D were
near-tied in F1 (0.568858 versus 0.565764). Configuration D achieved higher
precision and reduced the mean false-positive entity count from 584.0 to
529.0. These observations identify RAG as a precision- and
false-positive-oriented operating-point control that complements the best-F1
glossary configuration in this frozen low-label equipment-domain NER setting.
As a focused external check, we also
evaluated all 1,343 labeled
rows of the official CLUENER development split with the same A/B/C/D, $K=5$,
three-seed protocol; the official CLUENER test split is unlabeled and was not
used. Mean F1 for A--D was 0.428519, 0.513964, 0.517521, and 0.513424,
respectively. Thus, the directions for kNN, the support-derived schema-specific
lexicon, and the RAG precision/false-positive trade-off form a limited
descriptive replication on one additional Chinese NER benchmark; they do not
establish broad portability or cross-domain generalization. An accompanying
four-cell relation audit preserves the framework's exploratory
information-extraction boundary: the two-stage execution is reproducible and
auditable, but exact relation performance is low and every configuration has
zero \texttt{HAS\_METRIC} true positives. The supported performance
contribution is therefore the EC-FSIE NER route and its observed false-positive
trade-off.
